% --- Archon physics formalization source begin ---
% archon:physics
% archon:chemistry
% archon:covers IChO2026Problems/problem_icho_2026_t5_a1.lean
% archon:source-report reports/icho_2026/problem_icho_2026_t5_a1.source.json
% archon:problem-id icho_2026_t5
% archon:part-id T5-A1

\chapter{Icho Chemistry problem icho\_2026\_t5\_a1}
\label{ch:IChO2026Problems_problem_icho_2026_t5_a1}

\paragraph{Problem source.}
\#\# Source
58th International Chemistry Olympiad, Tashkent, Uzbekistan, 2026.
Official-materials index: https://www.icho2026.uz/problems
English problem PDF: ../icho\_2026\_source/raw/theory\_problem.pdf
English solution/rubric PDF: ../icho\_2026\_source/raw/theory\_solution.pdf
Problem PDF page 44; printed local question page 1; solution starts on PDF page 39.
The page PNG is primary visual evidence for formulas, structures, tables, and figures.

\#\# T5. Cardiolipins
T5. Cardiolipins 6\% In this problem you are asked to be creative and build molecules from fragments. For instance, if you are required to draw a chiral phospholipid such as W using fragments a-e, it can be assembled as shown below. Cardiolipins are a family of acyclic phospholipids differing in fatty acid residues. They are essential for energy production and membrane stability. Defects can lead to certain diseases. PL1 belongs to this family. Using only the structural elements a-d in the quantities stated below, it is possible to assemble the non-ionised form of PL1. R is a hydrocarbon substituent in the fatty acid structure. PL1 does not contain any peroxide bonds.

\#\# Current subquestion 5.1
Tick one correct statement regarding the value of n (number of type a fragments). (a) n is an even number, (b) n is an odd number, (c) n can be either an odd or an even number.

\#\# Dataset metadata
IChO 2026; T5; raw rubric points=1; kind=classification; formalization\_ready=true.

\paragraph{Shared problem context.}
T5. Cardiolipins 6\% In this problem you are asked to be creative and build molecules from fragments. For instance, if you are required to draw a chiral phospholipid such as W using fragments a-e, it can be assembled as shown below. Cardiolipins are a family of acyclic phospholipids differing in fatty acid residues. They are essential for energy production and membrane stability. Defects can lead to certain diseases. PL1 belongs to this family. Using only the structural elements a-d in the quantities stated below, it is possible to assemble the non-ionised form of PL1. R is a hydrocarbon substituent in the fatty acid structure. PL1 does not contain any peroxide bonds.

\paragraph{Current subquestion.}
Tick one correct statement regarding the value of n (number of type a fragments). (a) n is an even number, (b) n is an odd number, (c) n can be either an odd or an even number.

\paragraph{Recorded answer/context.}
Each bond involves two atoms, so the total number of atoms with broken bonds (indicated by the wavy line on the fragments of the molecule) must be an even number. For fragments with a known number of repeats in the PL1 molecule, the total number of atoms with broken bonds is 17 (2·2 + 3·3 + 4). Thus the total number of hydrogen atoms with broken bonds must be an odd number. Therefore, n is an odd number. Correct answer: 1 point. 1.0 pt

\paragraph{Figure/image path.}
../icho\_2026\_source/image/T5\_page-1.png

\paragraph{Formalization target.}
create a compiling Lean file with sorry bodies at `IChO2026Problems/problem\_icho\_2026\_t5\_a1.lean`.
The Lean declarations must preserve every mathematical object, quantifier, hypothesis, side condition, approximation/error guarantee, and requested conclusion from the source statement.
Use Mathlib/Physlib/Chemistry names found through LeanExplore where available. Prefer the configured benchmark Base library; introduce a local definition only when the source genuinely requires an object that the environment does not expose.

\begin{theorem}[Icho Chemistry formalization target]
\label{thm:physics:icho_2026_t5_a1:target}
The assigned autoformalize agent should translate this IChO chemistry problem into Lean declarations in the covered file, with theorem and lemma proof bodies written as `by sorry`.
\end{theorem}
\begin{proof}
This is an autoformalization task, not a proof task. Produce faithful statements that can later be proved without weakening the source contract.
\end{proof}
% --- Archon physics formalization source end ---
